\documentclass[11pt]{article}
\usepackage[T1]{fontenc}
\usepackage[utf8]{inputenc}
\usepackage{geometry}
\usepackage{hyperref}
\usepackage{xcolor}
\usepackage{minted}
\geometry{margin=1in}

\definecolor{xhbg}{HTML}{FAFAFA}
\setminted{
  fontsize=\small,
  breaklines=true,
  autogobble=true,
  frame=single,
  framesep=2mm,
  linenos=true,
  tabsize=2,
  bgcolor=xhbg
}
\usemintedstyle{friendly}

\newminted[code]{haskell}{}

\title{Xeus-Haskell: Docs and Literate Sources}
\author{}
\date{}

\begin{document}
\maketitle
\tableofcontents
\newpage

\section{Welcome to Xeus-Haskell}


Xeus-Haskell is a Jupyter kernel for Haskell based on the native implementation of the Jupyter protocol, \href{https://github.com/jupyter-xeus/xeus}{xeus}.
\subsection{Preliminaries}


To start using or developing Xeus-Haskell, please note the following:
\begin{itemize}
\item \textbf{MicroHs Backend:} Xeus-Haskell uses \href{https://github.com/augustss/MicroHs}{MicroHs}, a minimal implementation of Haskell. Therefore, supported libraries and extensions are limited.
\item \textbf{Native Protocol:} Based on \href{https://github.com/jupyter-xeus/xeus}{xeus} (C++), so the kernel must be compiled for the specific platform you are using.
\item \textbf{WebAssembly Support:} You can compile this kernel to WebAssembly to run it directly in browsers!
\end{itemize}


\subsection{Quickstart}


Since this project is not yet available on Conda-forge or Emscripten-forge, you need to build it from source. We have organized the workflows using \href{https://pixi.sh}{pixi}.
\subsubsection{For Native JupyterLab}


\begin{minted}{text}
git clone https://github.com/jupyter-xeus/xeus-haskell
cd xeus-haskell
pixi run -e default prebuild
pixi run -e default build
pixi run -e default install
pixi run -e default serve # JupyterLab is ready!
\end{minted}


\subsubsection{For WebAssembly JupyterLite}


\begin{minted}{text}
git clone https://github.com/jupyter-xeus/xeus-haskell
cd xeus-haskell
pixi install -e prebuild
pixi run -e wasm-build prebuild
pixi run -e wasm-build build
pixi run -e wasm-build install
# pixi run -e wasm-build fix-emscripten-links # Run this if you encounter linking errors
pixi run -e wasm-build serve # JupyterLite is ready!
\end{minted}


\subsection{A Deep Dive into Xeus-Haskell}


\begin{itemize}
\item \hyperref[sec:contributing]{Contributing Guide}
\item \hyperref[sec:communication]{Communication Protocol}
\item \hyperref[sec:repl]{REPL Mechanism}
\item \hyperref[sec:todo]{Roadmap \& TODO}
\item \href{https://github.com/jupyter-xeus/xeus-haskell/blob/main/AGENTS.md}{AI Agent Guide}
\end{itemize}

\section{Communication and System Architecture}
\label{sec:communication}


This document describes the internal communication flow of \texttt{xeus-haskell}, from the Jupyter frontend down to the MicroHs runtime.
\subsection{System Stack}


The communication stack is organized as follows:
Jupyter Frontend \(\leftrightarrow\) Xeus C++ Library \(\leftrightarrow\) Xeus-Haskell Kernel \(\leftrightarrow\) \texttt{mhs\_repl.cpp} Bridge \(\leftrightarrow\) Repl.lhs Compiler API \(\leftrightarrow\) MicroHs Runtime.


\subsection{Layer Details}


\subsubsection{1. Jupyter and Xeus Protocol Layer}


This layer handles the ZeroMQ-based Jupyter protocol. It manages the heartbeat, shell, and control sockets.
\subsubsection{2. Xeus to Kernel Interpreter Layer}


The \texttt{xinterpreter} implementation hooks into the Xeus kernel lifecycle. It receives execution requests and delegates them to the Haskell backend.
\subsubsection{3. C++ Bridge to Haskell Runtime Layer}


This is the bridge layer between C++ and the Haskell-compiled Compiler API wrapper (\texttt{Repl.lhs}).
The \texttt{mhs\_repl.cpp} bridge performs three primary tasks:
\begin{enumerate}
\item \textbf{Initialize MicroHs Runtime}:
   The \texttt{MicroHsRepl} constructor locates runtime files (standard libraries) using \texttt{microhs\_runtime\_dir}. After initialization, it evaluates a trivial expression (\texttt{"0"}) to warm up the compiler and cache library compilation.
\item \textbf{Capture Standard Output}:
   To receive output from the Haskell world, we use POSIX pipes to redirect \texttt{stdout}.
   \begin{itemize}
   \item C++ sets up a redirection from \texttt{stdout} to a temporary buffer.
   \item After Haskell evaluation completes, the buffer is read.
   \item For implementation details, see \texttt{capture\_stdout} in the source code.
   \end{itemize}
\item \textbf{Parse Captured Content}:
   The output is parsed to distinguish plain text from rich content (HTML, LaTeX, etc.) using the internal protocol described below.
\end{enumerate}


\subsubsection{4. REPL Engine and MicroHs Compiler Layer}


This layer uses the MicroHs Compiler API to evaluate Haskell code strings in the runtime environment.
\subsection{Rich Content Protocol}


\texttt{xeus-haskell} supports rendering rich media (HTML, images, LaTeX) by embedding control sequences in the standard output stream.
\subsubsection{Protocol Structure}


Rich content is strictly delimited by \textbf{ASCII control characters}:
\begin{minted}{text}
<STX><MIME_TYPE><US><CONTENT><ETX>
\end{minted}


\begin{center}
\begin{tabular}{|l|c|p{8cm}|}
\hline
\textbf{Component} & \textbf{ASCII (Hex)} & \textbf{Description} \\
\hline
\texttt{STX} & \texttt{02} & \textbf{S}tart of \textbf{T}e\textbf{x}t \\
\hline
\texttt{MIME\_TYPE} & - & The target MIME type (e.g., \texttt{text/html}) \\
\hline
\texttt{US} & \texttt{1F} & \textbf{U}nit \textbf{S}eparator \\
\hline
\texttt{CONTENT} & - & The raw data/string to be rendered \\
\hline
\texttt{ETX} & \texttt{03} & \textbf{E}nd of \textbf{T}e\textbf{x}t \\
\hline
\end{tabular}
\end{center}
\subsubsection{Supported Content Types}


\begin{itemize}
\item \texttt{text/plain}: Standard console output.
\item \texttt{text/html}: Rendered as HTML in Jupyter.
\item \texttt{text/latex}: Mathematical equations rendered via MathJax.
\item \texttt{text/markdown}: Formatted text and tables.
\end{itemize}


All content is expected to be \textbf{UTF-8} encoded.
\subsection{Haskell Display API}


The protocol is abstracted in Haskell through the \texttt{Display} typeclass.
\subsubsection{1. Display Typeclass Interface}


To enable rich output for a custom type, implement the \texttt{Display} typeclass:
\begin{minted}{haskell}
import XHaskell.Display

-- | Data structure for the protocol
data DisplayData = DisplayData {
    mimeType :: String,
    content  :: String
}
\end{minted}


\subsubsection{2. Implementation Examples}


\begin{minted}{haskell}
newtype HTMLString = HTMLString String
instance Display HTMLString where
    display (HTMLString s) = DisplayData "text/html" s

newtype LaTeXString = LaTeXString String
instance Display LaTeXString where
    display (LaTeXString s) = DisplayData "text/latex" s

-- Default instance for plain strings
instance Display String where
    display s = DisplayData "text/plain" s
\end{minted}


\subsubsection{3. Protocol Generation}


The \texttt{DisplayData} record implements a \texttt{Show} instance that automatically formats the string into the \texttt{<STX>...<ETX>} protocol.
\begin{minted}{haskell}
putStr $ show (display (HTMLString "<p style='color: red'>Hello Haskell!</p>"))
\end{minted}


\texttt{02} \texttt{text/html} \texttt{1F} \texttt{<p style='color: red'>Hello Haskell!</p>} \texttt{03}

\section{MicroHs REPL Mechanism}
\label{sec:repl}

This chapter documents the REPL as a single argument: a Jupyter kernel can remain interactive and semantically coherent even when its compiler was not designed as a notebook runtime. The architectural choice of MicroHs is deliberate. A small compiler surface constrains the language to a pragmatic Haskell 2010 core, but it also exposes a tractable runtime boundary for foreign-function integration, deployment to constrained environments, and careful control of evaluation latency.

From a scientific perspective, the system is a state machine with explicit transitions and measurable costs. From a humanistic perspective, it is a continuity device: each cell is a new utterance in an ongoing discourse, and the runtime must preserve memory, revision, and context without losing formal precision.

\subsection{Layered Architecture: Notebook Request to Runtime Effect}

Execution is a composition of layers
$L_0 \to L_1 \to L_2 \to L_3 \to L_4$:

\begin{itemize}
\item $L_0$: Jupyter message protocol (execute, complete, inspect, is\_complete).
\item $L_1$: Xeus interpreter callbacks in C++.
\item $L_2$: C++ bridge that captures output and marshals FFI buffers.
\item $L_3$: Haskell FFI entrypoints managing persistent \verb|ReplCtx|.
\item $L_4$: MicroHs parse/typecheck/translate/evaluate pipeline.
\end{itemize}

The Xeus callback layer delegates kernel behavior to the bridge object:

\begin{minted}{cpp}
void interpreter::execute_request_impl(send_reply_callback cb,
                                       int execution_counter,
                                       const std::string &code,
                                       xeus::execute_request_config config,
                                       nl::json /*user_expressions*/) {
  auto exec_result = [&]() -> repl_result {
    try {
      return m_repl.execute(code);
    } catch (const std::exception &e) {
      return {false, std::string(), std::string(e.what())};
    } catch (...) {
      return {false, std::string(), std::string("Unknown MicroHs error")};
    }
  }();

  if (!exec_result.ok) {
    const std::string &error_msg = exec_result.error;
    const std::vector<std::string> traceback{error_msg};
    publish_execution_error("RuntimeError", error_msg, traceback);

    nl::json traceback_json = nl::json::array();
    traceback_json.push_back(error_msg);

    cb(xeus::create_error_reply(error_msg, "RuntimeError", traceback_json));
    return;
  }

  if (!config.silent) {
    const std::string &output = exec_result.output;
    if (!output.empty()) {
      nl::json pub_data;
      pub_data[exec_result.mime_type] = output;
      publish_execution_result(execution_counter, std::move(pub_data),
                               nl::json::object());
    }
  }

  cb(xeus::create_successful_reply(nl::json::array(), nl::json::object()));
}

nl::json interpreter::is_complete_request_impl(const std::string &code) {
  std::string status = m_repl.is_complete(code);
  return xeus::create_is_complete_reply(status /*status*/, "   " /*indent*/);
}
\end{minted}

The bridge constructor initializes MicroHs once, discovers runtime paths, and triggers a warmup evaluation:

\begin{minted}{cpp}
MicroHsRepl::MicroHsRepl() {
  if (!mhs_repl_initialized) {
    mhs_init();
    mhs_repl_initialized = true;
  }

  std::string microhs_runtime_dir;
  const auto* mhsdir = std::getenv("MHSDIR");
  const auto* conda_prefix = std::getenv("CONDA_PREFIX");

  if (mhsdir != nullptr) {
    microhs_runtime_dir = mhsdir;
  } else if (conda_prefix != nullptr) {
    microhs_runtime_dir = std::string(conda_prefix) + "/share/microhs";
  }

  context = mhs_repl_new(const_cast<char*>(microhs_runtime_dir.c_str()),
                         static_cast<uintptr_t>(microhs_runtime_dir.size()));

  auto warmup = execute("0");
  (void)warmup;
}

repl_result MicroHsRepl::execute(std::string_view code) {
  std::string raw_output;

  try {
    raw_output = capture_stdout([&]() {
      std::string code_str(code);
      char* err = nullptr;
      intptr_t rc =
          mhs_repl_execute(context, const_cast<char*>(code_str.c_str()),
                           static_cast<uintptr_t>(code_str.size()),
                           reinterpret_cast<void*>(&err));

      if (rc != 0 && err) {
        std::string err_str(err);
        mhs_repl_free_cstr(err);
        throw std::runtime_error(err_str);
      }
      if (err) {
        mhs_repl_free_cstr(err);
      }
    });
  } catch (const std::runtime_error& e) {
    return {false, std::string(), std::string(e.what()), "text/plain"};
  }

  raw_output.erase(std::remove(raw_output.begin(), raw_output.end(), '\r'),
                   raw_output.end());

  ParsedOutput parsed = parse_protocol_output(raw_output);

  return {true, std::move(parsed.content), std::string(),
          std::move(parsed.mime_type)};
}
\end{minted}

On the Haskell boundary, the FFI adapter centralizes exception normalization, status-code mapping, and state updates:

\begin{minted}{haskell}
runStatefulAction
  :: (ReplCtx -> String -> IO (Either ReplError ReplCtx))
  -> ReplHandle
  -> CString
  -> CSize
  -> Ptr CString
  -> IO CInt
runStatefulAction act h srcPtr srcLen errPtr =
  withHandleInput h srcPtr srcLen $ \ref src -> do
    ctx <- readIORef ref
    result <- try (act ctx src) :: IO (Either SomeException (Either ReplError ReplCtx))
    case normalizeResult result of
      Left err -> writeErrorCString errPtr (prettyReplError err) >> pure c_ERR
      Right ctx' -> writeIORef ref ctx' >> pure c_OK

mhsReplExecute :: ReplHandle -> CString -> CSize -> Ptr CString -> IO CInt
mhsReplExecute = runStatefulAction replExecute

mhsReplIsComplete :: ReplHandle -> CString -> CSize -> IO CString
mhsReplIsComplete h srcPtr srcLen =
  withHandleInput h srcPtr srcLen $ \ref src -> do
    ctx <- readIORef ref
    status <- replIsComplete ctx src
    newCString status
\end{minted}

\subsection{Persistent Context as a Formal Object}

The session state is represented by:
$$
ReplCtx = (F, C, D, \Sigma)
$$
where $F$ is compiler flags/search paths, $C$ is MicroHs incremental cache, $D$ is ordered stored definitions, and $\Sigma$ is the pair of type/value symbol tables.

This is reflected directly in the source:

\begin{minted}{haskell}
type ReplHandle = StablePtr (IORef ReplCtx)

data StoredDef = StoredDef
  { sdCode :: String
  , sdNames :: [Ident]
  }

data ReplCtx = ReplCtx
  { rcFlags :: Flags
  , rcCache :: Cache
  , rcDefs  :: [StoredDef]
  , rcSyms  :: Symbols
  }

initialCtx :: String -> IO ReplCtx
initialCtx dir = do
  let flags = defaultFlags dir
  mpath <- lookupEnv "MHS_LIBRARY_PATH"
  let rpath = maybe "." id mpath
      extra = splitColon rpath
      rcFlags = flags { paths = paths flags ++ extra }
  pure ReplCtx { rcFlags = rcFlags, rcCache = emptyCache, rcDefs = [], rcSyms = (stEmpty, stEmpty) }
\end{minted}

Definition persistence uses replacement-by-name to model REPL shadowing:

\begin{minted}{haskell}
appendDefinition :: ReplCtx -> String -> Either ReplError [StoredDef]
appendDefinition ctx snippet =
    case extractDefinitionNames snippet of
        Left err -> Left err
        Right names ->
            let uniqueNames = nub names
                retainedDefs = stripRedefined (rcDefs ctx) uniqueNames
            in Right (retainedDefs ++ [StoredDef snippet uniqueNames])
\end{minted}

Formally, with $\mathrm{Names}(d)$ as bound identifiers in definition snippet $d$,
$$
D' = \mathrm{strip}(D, \mathrm{Names}(d)) \mathbin{\|} [d].
$$
The invariant is chronological except where name-collision replacement occurs; the most recent binding for each identifier is retained.

\subsection{Scan-and-Split as a Parsing Strategy}

Notebook cells may mix declarations and runnable expressions. For lines
$\ell = \langle \ell_1, \dots, \ell_n \rangle$, the kernel searches split index $i \in \{n, n-1, \dots, 0\}$ and partitions:
$$
\ell^{def}_i = \langle \ell_1, \dots, \ell_i \rangle,
\qquad
\ell^{run}_i = \langle \ell_{i+1}, \dots, \ell_n \rangle.
$$
A valid plan is the first $i$ such that accumulated definitions parse and the tail is either empty, expression-valid, or expression-valid after \verb|do|-wrapping.

\begin{minted}{haskell}
data SplitPlan
  = SplitDefineOnly String
  | SplitDefineThenRun String String

firstValidSplitPlan :: String -> String -> Maybe SplitPlan
firstValidSplitPlan currentDefs snippet =
  let snippetLines = lines (ensureTrailingNewline snippet)
  in listToMaybe (mapMaybe (classifySplit currentDefs snippetLines) [length snippetLines, length snippetLines - 1 .. 0])

classifySplit :: String -> [String] -> Int -> Maybe SplitPlan
classifySplit currentDefs snippetLines splitIndex
  | not (canParseDefinition candidateDefs) = Nothing
  | all allwsLine runLines = Just (SplitDefineOnly defPart)
  | canParseExpression runPart = Just (SplitDefineThenRun defPart runPart)
  | canParseExpression doRunPart = Just (SplitDefineThenRun defPart doRunPart)
  | otherwise = Nothing
  where
    (defLines, runLines) = splitAt splitIndex snippetLines
    defPart = unlines defLines
    runPart = unlines (dropWhileEndLocal allwsLine runLines)
    doRunPart = "do\n" ++ indent runPart
    candidateDefs = currentDefs ++ defPart
\end{minted}

The completeness classifier reuses this split logic and adds structural heuristics:

\begin{minted}{haskell}
isIncomplete :: String -> Bool
isIncomplete s = go [] s
  where
    go st [] = not (null st)
    go st (c:cs)
      | c `elem` "([{" = go (c:st) cs
      | c `elem` ")]}" = case st of
          [] -> False
          (x:xs) -> if match x c then go xs cs else False
      | c == '"' = goString st cs
      | c == '\'' = goChar st cs
      | otherwise = go st cs
    goString st [] = True
    goString st ('\\':'"':cs) = goString st cs
    goString st ('"':cs) = go st cs
    goString st (_:cs) = goString st cs
    goChar st [] = True
    goChar st ('\\':'\'':cs) = goChar st cs
    goChar st ('\'':cs) = go st cs
    goChar st (_:cs) = goChar st cs
    match '(' ')' = True
    match '[' ']' = True
    match '{' '}' = True
    match _ _ = False
\end{minted}

The returned frontend status can be expressed as:
$$
\exists i\,\mathrm{validSplit}(x,i) \Rightarrow \mathcal{C}(x)=\mathsf{complete}
$$
$$
\neg\exists i\,\mathrm{validSplit}(x,i) \land \mathrm{openDelims}(x)
\Rightarrow \mathcal{C}(x)=\mathsf{incomplete}
$$
$$
\neg\exists i\,\mathrm{validSplit}(x,i) \land \neg\mathrm{openDelims}(x)
\Rightarrow \mathcal{C}(x)=\mathsf{invalid}
$$

\subsection{Compilation and Evaluation Transitions}

Code generation and execution are isolated into a small compilation module:

\begin{minted}{haskell}
runBlock :: String -> String
runBlock stmt = unlines
  [ runResultName ++ " :: IO ()"
  , runResultName ++ " = _printOrRun ("
  , indent stmt
  , "  )"
  ]

compileModule :: ReplCtx -> String -> IO (Either ReplError (TModule [LDef], Cache, Symbols))
compileModule ctx src =
  case parse pTopModule "<repl>" src of
    Left perr -> pure (Left (ReplParseError perr))
    Right mdl -> do
      r <- runStateIO (compileModuleP (rcFlags ctx) ImpNormal mdl) (rcCache ctx)
      let (((dmdl, syms, _, _, _), _), cache') = unsafeCoerce r
          cmdl = compileToCombinators dmdl
      pure (Right (cmdl, cache', syms))

runAction :: Cache -> TModule [LDef] -> Ident -> IO (Either ReplError ())
runAction cache cmdl ident = do
  let defs      = tBindingsOf cmdl
      baseMap   = withBuiltinAliases $ translateMap $ concatMap tBindingsOf (cachedModules cache)
      actionAny = translateWithMap baseMap (defs, Var (qualIdent (tModuleName cmdl) ident))
      action    = unsafeCoerce actionAny :: IO ()
  action
  pure (Right ())
\end{minted}

At the executor level, the transitions are explicit:

\begin{minted}{haskell}
replDefine :: ReplCtx -> String -> IO (Either ReplError ReplCtx)
replDefine ctx snippet =
  case appendDefinition ctx (ensureTrailingNewline snippet) of
    Left err -> pure (Left err)
    Right defsWithNew ->
      compileModule ctx (moduleFromDefs defsWithNew) `andThen`
        \(_, cache', syms') ->
          right ctx{ rcDefs = defsWithNew, rcCache = cache', rcSyms = syms' }

replRun :: ReplCtx -> String -> IO (Either ReplError ReplCtx)
replRun ctx stmt =
  compileModule ctx (moduleSourceWith ctx (runBlock stmt)) `andThen`
    \(cmdl, cache', syms') ->
      runAction cache' cmdl runResultIdent `andThen`
        \() -> right ctx{ rcCache = cache', rcSyms = syms' }

replExecute :: ReplCtx -> String -> IO (Either ReplError ReplCtx)
replExecute ctx snippet
  | Just cmd <- parseMetaCommand snippet = runMetaCommand ctx cmd
replExecute ctx snippet =
  case firstValidSplitPlan (currentDefsSource ctx) snippet of
    Nothing -> pure (Left (ReplParseError "unable to parse snippet"))
    Just (SplitDefineOnly defPart) -> replDefine ctx defPart
    Just (SplitDefineThenRun defPart runPart) -> defineThenRun ctx defPart runPart

defineThenRun :: ReplCtx -> String -> String -> IO (Either ReplError ReplCtx)
defineThenRun ctx defPart runPart = replDefine ctx defPart `andThen` \ctx' -> replRun ctx' runPart

replIsComplete :: ReplCtx -> String -> IO String
replIsComplete ctx snippet = pure (toText completion)
  where
    completion
      | all isws snippet = Complete
      | hasValidSplit = Complete
      | isIncomplete snippet = Incomplete
      | otherwise = Invalid
    hasValidSplit = isJust (firstValidSplitPlan (currentDefsSource ctx) snippet)
    toText state =
      case state of
        Complete -> "complete"
        Incomplete -> "incomplete"
        Invalid -> "invalid"
\end{minted}

We can summarize the dynamic semantics as two primary transitions:
$$
\delta_{def}(F,C,D,\Sigma; d) = (F, C', D', \Sigma') \quad \mathsf{or} \quad \mathsf{error},
$$
$$
\delta_{run}(F,C,D,\Sigma; e) = (F, C', D, \Sigma') \quad \mathsf{or} \quad \mathsf{error}.
$$
The definitional transition mutates $D$; the run transition preserves $D$ but still advances cache and symbols.

\subsection{Error Algebra and Utility Substrate}

The error type is intentionally small and compositional:

\begin{minted}{haskell}
data ReplError
  = ReplParseError String
  | ReplCompileError String
  | ReplRuntimeError String
  deriving (Eq, Show)

prettyReplError :: ReplError -> String
prettyReplError e =
  case e of
    ReplParseError s   -> "Parse error: "   ++ s
    ReplCompileError s -> "Compile error: " ++ s
    ReplRuntimeError s -> "Runtime error: " ++ s
\end{minted}

The support module codifies lexical whitespace and synthetic module construction used throughout analysis and execution:

\begin{minted}{haskell}
ensureTrailingNewline :: String -> String
ensureTrailingNewline s
  | null s         = "\n"
  | last s == '\n' = s
  | otherwise      = s ++ "\n"

moduleHeader :: String
moduleHeader = unlines
  [ "module Inline where"
  , "import Prelude"
  , "import System.IO.PrintOrRun"
  , "import Data.Typeable"
  , "import Numeric"
  ]

buildModule :: String -> String
buildModule defs = moduleHeader ++ defs
\end{minted}

These helpers are modest in size, but conceptually central: they preserve the textual discipline that allows parser checks, split classification, and compile invocation to agree on what a ``module'' means.

\subsection{Runtime Constraints and Caching Strategy}

MicroHs imposes a productive constraint: the kernel targets a portable subset centered on Haskell 2010 behavior. The tradeoff is intentional. Instead of pursuing maximal language surface, the implementation optimizes for reproducible compilation paths, tractable FFI contracts, and startup predictability in local and WebAssembly contexts.

Latency can be viewed as:
$$
T_{first} = T_{init} + T_{runtime\_path} + T_{prelude\_compile} + T_{cell},
$$
$$
T_{next} \approx T_{cache\_reuse} + T_{delta\_compile} + T_{cell}.
$$
The warmup evaluation in the bridge lowers the first user-visible cell by moving a portion of $T_{prelude\_compile}$ into initialization. The persistent cache $C$ inside \verb|ReplCtx| then amortizes subsequent compilation by reusing previously compiled modules and symbol context.

The architecture-level implication is that ``state'' is not only user definitions; it is also compilation history. In notebook practice, this turns a sequence of independent source texts into a cumulative dialogue with memory.

\subsection{Relation to Rich Output Protocol}

This chapter has focused on control flow and semantic state. The rendering path for MIME-aware outputs is documented in \hyperref[sec:communication]{Communication and System Architecture}. Operationally, the REPL emits values and effects; the communication protocol decides how those effects are materialized for the frontend.

\section{Contributing to xeus-haskell}
\label{sec:contributing}


Thank you for your interest in contributing to \texttt{xeus-haskell}!
Xeus and xeus-haskell are subprojects of Project Jupyter and are subject to the \href{https://github.com/jupyter/governance}{Jupyter governance} and \href{https://github.com/jupyter/governance/blob/master/conduct/code\_of\_conduct.md}{Code of conduct}.
\subsection{General Guidelines}


For general documentation about contributing to Jupyter projects, see the \href{https://jupyter.readthedocs.io/en/latest/contributor/content-contributor.html}{Project Jupyter Contributor Documentation}.
\subsection{Community}


The Xeus team organizes public video meetings. The schedule for future meetings and minutes of past meetings can be found on our \href{https://jupyter-xeus.github.io/}{Team Compass}.
\subsection{Getting Started}


\subsubsection{Prerequisites}


The project uses \href{https://pixi.sh/}{Pixi} for dependency management and workflow automation. If you haven't installed it yet:
\begin{minted}{text}
curl -fsSL https://pixi.sh/install.sh | sh
\end{minted}


\subsubsection{Setting Up Your Environment}


\begin{enumerate}
\item Fork the repository on GitHub.
\item Clone your fork locally:
\end{enumerate}


\begin{minted}{text}
git clone https://github.com/<your-username>/xeus-haskell.git
cd xeus-haskell
\end{minted}


\begin{enumerate}
\item Initialize the development environment:
\end{enumerate}


\begin{minted}{text}
pixi run -e default prebuild
\end{minted}


The \texttt{default} environment includes all necessary tools (CMake, C++ compilers, Python for testing, etc.).
\subsection{Development Workflow}


\subsubsection{Building the Kernel}


To build and install the kernel in your local environment:
\begin{minted}{text}
pixi run -e default build
pixi run -e default install
\end{minted}


\subsubsection{Running Tests}


We use \texttt{pytest} and \texttt{ctest} for verifying the kernel:
\begin{minted}{text}
# Run C++ tests
pixi run -e default ctest

# Run Python-based Jupyter kernel tests
pixi run -e default pytest
\end{minted}


\subsubsection{Working with WebAssembly (JupyterLite)}


If you are contributing to the WebAssembly build, you need the \texttt{wasm-build} and \texttt{wasm-host} environments:
\begin{minted}{text}
# Prepare the WASM host environment
pixi install -e wasm-build
pixi install -e wasm-host

# Build for WebAssembly
pixi run -e wasm-build prebuild
pixi run -e wasm-build build
pixi run -e wasm-build install

# Run the JupyterLite server locally
pixi run -e wasm-build serve
\end{minted}


\subsection{Project Structure}


\begin{itemize}
\item \texttt{src/}: C++ source code for the kernel and REPL bridge.
\item \texttt{include/}: C++ header files.
\item \texttt{share/}: Haskell side of the implementation, including \texttt{Repl.lhs} and MicroHs libraries.
\item \texttt{test/}: Python tests for kernel functionality.
\item \texttt{xhaskell.tex}: Integrated LaTeX manual source that includes all architecture chapters and literate Haskell modules.
\end{itemize}


\subsection{Style Guidelines}


\begin{itemize}
\item \textbf{C++}: Follow modern C++ practices. We use Xeus as the base library.
\item \textbf{Haskell}: Since we use MicroHs, avoid dependencies or language features not supported by the MicroHs compiler.
\item \textbf{Documentation}: Write documentation in LaTeX (\texttt{.tex}) and keep examples runnable with the current Pixi workflows.
\end{itemize}

\subsubsection{Building Documentation}

Generate the project PDF documentation with:
\begin{minted}{text}
pixi run -e docs generate
\end{minted}


\subsection{Submitting Changes}


\begin{enumerate}
\item Create a new branch for your feature or bugfix.
\item Ensure all tests pass.
\item Submit a Pull Request with a clear description of the changes.
\end{enumerate}


\subsection{Community}


Join our public video meetings! The schedule and minutes are available on our \href{https://jupyter-xeus.github.io/}{Team Compass}.

\section{TODO: Unimplemented Features \& Roadmap}
\label{sec:todo}


This document tracks planned improvements and currently missing features for \texttt{xeus-haskell}.
\subsection{Core Kernel Features}


\begin{itemize}
\item [x] \textbf{\texttt{is\_complete} implementation}: Map parsability of mixed cells to Jupyter's completeness status.
\item [ ] \textbf{Incomplete Status}: Refine \texttt{is\_complete} to distinguish between \texttt{invalid} and \texttt{incomplete} code (e.g., unclosed delimiters).
\item [x] \textbf{\texttt{inspect\_request}}: Support for "Introspection" (Shift+Tab in Jupyter) to show documentation or type signatures for identifiers.
\item [ ] \textbf{\texttt{history\_request}}: Implementation of the Jupyter history protocol to allow searching and retrieving previous cell inputs.
\item [ ] \textbf{Advanced Completion}: Improve \texttt{completion\_request} to support qualified names (e.g., \texttt{Prelude.putStrLn}) and type-aware suggestions.
\item [ ] \textbf{Kernel Interrupt}: Add support for interrupting long-running Haskell executions (SIGINT handling).
\end{itemize}


\subsection{REPL Improvements}


\begin{itemize}
\item [ ] \textbf{Selective Re-compilation}: Instead of concatenating all previous definitions for every new definition, implement a dependency-tracking system to only re-compile what is necessary.
\item [ ] \textbf{Better Error Reporting}: Map MicroHs compile errors back to the specific line numbers in the Jupyter cell.
\item [ ] \textbf{WASM Optimization}: Reduce the size of the WebAssembly bundle and optimize the initial "warmup" time in the browser.
\item [x] \textbf{Language Extensions}: Work with the MicroHs upstream to support more modern Haskell extensions (e.g., \texttt{GADTs}, \texttt{TypeFamilies}).
\end{itemize}


\subsection{Rich Output Display System}


\begin{itemize}
\item [ ] \textbf{Standard Library Instances}: Add \texttt{Display} instances for more types in the MicroHs standard library.
\item [ ] \textbf{Data Visualization}: Create bridges for common Haskell charting libraries to output Jupyter-compatible JSON/Vega-Lite data.
\item [ ] \textbf{Widgets}: Initial investigation into supporting Jupyter Widgets (ipywidgets protocol).
\end{itemize}


\subsection{Infrastructure \& Testing}


\begin{itemize}
\item [ ] \textbf{Enhanced CI}: Add automated UI tests using \texttt{playwright} or \texttt{selenium} to verify the kernel works in a real JupyterLab/JupyterLite instance.
\item [ ] \textbf{Benchmarking}: Create a suite of performance benchmarks to track compilation and execution speed regressions.
\item [ ] \textbf{Conda-forge/Emscripten-forge}: Package the kernel for easier installation without needing to build from source.
\end{itemize}

\clearpage
\section{Literate Haskell Source}
\subsection{REPL Foreign-Function Interface Module}
\input{src/Repl.lhs}
\subsection{Snippet Analysis and Parsing Module}
\input{src/Repl/Analysis.lhs}
\subsection{Compilation and Execution Translation Module}
\input{src/Repl/Compiler.lhs}
\subsection{Persistent Context and Definition Storage Module}
\input{src/Repl/Context.lhs}
\subsection{Error Classification and Rendering Module}
\input{src/Repl/Error.lhs}
\subsection{Runtime Command Execution Module}
\input{src/Repl/Executor.lhs}
\subsection{Shared Utility Functions Module}
\input{src/Repl/Utils.lhs}

\end{document}
